% Generated draft: requires mathematical and source review.
\documentclass[11pt]{article}
\usepackage[margin=25mm]{geometry}
\usepackage{fontspec}
\setmainfont{Latin Modern Roman}
\usepackage{amsmath,amssymb,mathtools,unicode-math}
\setmathfont{Latin Modern Math}
\usepackage{enumitem,needspace,xurl,hyperref}
\setlength{\parindent}{0pt}
\setlength{\parskip}{5pt}
\setlength{\emergencystretch}{3em}
\newcommand{\N}{\mathbb{N}}
\newcommand{\Z}{\mathbb{Z}}
\newcommand{\Q}{\mathbb{Q}}
\newcommand{\R}{\mathbb{R}}
\newcommand{\C}{\mathbb{C}}
\newcommand{\F}{\mathbb{F}}
\newcommand{\A}{\mathbb{A}}
\newcommand{\PP}{\mathbb{P}}
\newcommand{\E}{\mathbb{E}}
\newcommand{\eps}{\varepsilon}
\newcommand{\dd}{\,\mathrm{d}}
\begin{document}
\section*{UnsolvedMath ID: 701458}
\textbf{Source problem code:} catalogue entry
\subsection{Definitions and mathematical statement}
All groups considered are finite. A group is soluble if its derived series terminates at the trivial group. A Fitting class \(\mathfrak F\) of finite soluble groups is a nonempty class, closed under isomorphisms, such that every normal subgroup of a member of \(\mathfrak F\) belongs to \(\mathfrak F\), and whenever \(N_1,N_2\trianglelefteq G\) belong to \(\mathfrak F\), with \(G\) finite soluble, their product \(N_1N_2\) belongs to \(\mathfrak F\).

For \(H\leq G\), write \(H^g=g^{-1}Hg\). The subgroup \(H\) is abnormal in \(G\) if
\[
 g\in\langle H,H^g\rangle\qquad\text{for every }g\in G.
\]
An \(\mathfrak F\)-subgroup means a subgroup whose underlying group belongs to \(\mathfrak F\). For subgroups \(A,B\leq G\), the equation \(G=AB\) means that every element of \(G\) is a product \(ab\), where \(a\in A\) and \(b\in B\).

A formation of finite soluble groups is a nonempty isomorphism-closed class closed under homomorphic images and finite subdirect products. A subdirect product of \(G_1,\ldots,G_r\), where \(r\geq1\), is a subgroup of \(G_1\times\cdots\times G_r\) whose projection onto each factor is surjective.

\textbf{Question (A. F. Vasil'ev, Kourovka Notebook Problem 18.29).} Let \(\mathfrak F\) be a Fitting class of finite soluble groups. Assume that for every finite soluble group \(G\) and every pair of subgroups \(A,B\leq G\),
\[
\bigl(G=AB,\quad A,B\in\mathfrak F,\quad A\text{ and }B\text{ abnormal in }G\bigr)
\ \Longrightarrow\ G\in\mathfrak F.
\]
Must \(\mathfrak F\) be a formation?

The phrase ``abnormal \(\mathfrak F\)-subgroups'' is read here as subgroups belonging to \(\mathfrak F\) that are abnormal in the ambient group \(G\). The supplied record gives no separate definition of this phrase.
\subsection{Short English statement}
Suppose a Fitting class of finite soluble groups contains every finite soluble group that is a product of two abnormal subgroups belonging to the class. Must the class also be closed under homomorphic images and finite subdirect products?
\subsection{Sources}
\begin{itemize}
\item Local UnsolvedMath registry: supplied complete record with permanent ID \(701458\), source code \texttt{catalogue entry}, title \emph{Kourovka Notebook Problem 18.29}, and upstream revision \texttt{d371ee917ac9647fbb06cd6fc529deec6f3dda82}.
\item Supplied source attribution: \emph{The Kourovka Notebook (Unsolved Problems in Group Theory)}, 18th issue (2014), Problem 18.29; also identified in the record with the 21st issue edition (2026), \texttt{arXiv:1401.0300v48}. Proposed by A. F. Vasil'ev in 2014. The supplied collection description names E. I. Khukhro and V. D. Mazurov as editors.
\item The local record labels the problem \texttt{open}. This is supplied registry metadata, not an independently verified assessment of the literature. The cited source has not been independently consulted for this draft.
\end{itemize}
\end{document}
